% 08 · PLSCAN — persistence instead of a size floor
\section{PLSCAN: DROPPING THE SIZE FLOOR}
\label{sec:plscan}

\deckslides{slides 27--31 --- peaks, persistence, and the bars that survive.}

HDBSCAN$^*$ removed DBSCAN's density threshold but kept a single global
\texttt{min\_cluster\_size}: a point count, below which a structure is treated
as noise. That floor is a bad fit for our sample, where clusters range from 6
to 230 members. PLSCAN --- persistent multiscale density-based clustering,
\citet{Bot:25} --- takes the last step and removes it, replacing ``is this
group big enough?'' with ``how long does this group persist?''.

\subsection{Scale-space, and what persistence measures}

The idea comes from topological data analysis. Rather than committing to one
density threshold, consider the whole family of density estimates obtained by
smoothing the data at every scale. A real cluster appears as a \emph{peak}
that stays separate from its neighbours as the scale changes; a fluctuation
appears as a peak that is absorbed almost immediately. The range of scales
over which a peak exists as an independent object is its \emph{persistence},
and it is a statement about the data's structure rather than about a count.

The algorithm builds this structure explicitly:

\begin{enumerate}
\item \textbf{Scale-space.} Estimate the density field over a range of scales
  (equivalently: over a range of effective minimum cluster sizes), tracking the
  peaks at each scale.
\item \textbf{Leaf tree.} Record the peaks and their merges as a tree, where a
  \emph{leaf} is a density peak that exists in its own right --- the analogue
  of the branches of the HDBSCAN$^*$ condensed tree, but indexed by scale
  rather than by $k$-th-neighbour distance.
\item \textbf{Persistence trace.} For each leaf, measure the interval of scale
  over which the cluster survives as a separate object. Clusters by
  construction cannot be shorter-lived than their own persistence, which is
  what makes the criterion self-consistent.
\item \textbf{Selection.} Rate every cut of the leaf tree by the total
  persistence of the leaves it keeps, and take the cut with the highest total
  persistence; read off the labels at that cut.
\end{enumerate}

\begin{marginnote}[]
\entry{Persistence}{The range of minimum cluster sizes over which a cluster
exists as a separate structure, $s_{\max} - s_{\min}$ \citep{Bot:25}. A
long-lived peak is structure; a short-lived one is noise. It is measured in
stars --- the same units as the size floor it replaces --- which is why a
6-member cluster can outlive a 60-member fluctuation.}
\entry{Barcode}{One horizontal bar per cluster, its length equal to its
persistence, drawn on a common scale axis. Reading the barcode is the whole
selection step.}
\end{marginnote}

\begin{figure}[htbp]\centering
  \fig[0.94]{plscan_film.png}
  \caption{The persistence sweep, four frames of the deck's animation: the
  minimum-cluster-size axis is swept, and each peak's bar grows for as long as
  that peak remains a distinct cluster. Bars that stop early are fluctuations;
  bars that run the length of the sweep are structure.}
  \label{fig:plscanfim}
\end{figure}

\begin{figure}[htbp]\centering
  \begin{tabular}{@{}c@{\hskip5pt}c@{\hskip5pt}c@{}}
    \fig[0.30]{plscan_step1.png} &
    \fig[0.30]{plscan_step2.png} &
    \fig[0.30]{plscan_step3.png}
  \end{tabular}
  \caption{The three moves: \panel{a} find the density peaks; \panel{b} measure
  each peak's persistence and draw the barcode; \panel{c} keep the bars of the
  cut with the largest total persistence. The barcode replaces the size floor
  with a quantity measured over a range of sizes; the neighbour count $k$ that
  defines the core distance is still a choice, and when the trace has several
  peaks, so is which peak to report.}
  \label{fig:plscansteps}
\end{figure}

\begin{figure}[htbp]\centering
  \fig[0.78]{plscan_barcode.png}
  \caption{A three-peak density profile and its persistence barcode. The
  vertical extent of each bar is the range of ``minimum cluster size'' over
  which the corresponding cluster exists; the horizontal position identifies
  the peak. A min-cluster-size floor would cut all three bars at the same
  height, which is exactly the assumption PLSCAN avoids.}
  \label{fig:plscanbarcode}
\end{figure}

\subsection{What changes, and what does not}

Two honest statements, because the summary of this method is often oversold.

\textbf{What changes.} The criterion for ``is this a cluster?'' is no longer
a single size. A 6-member cluster that is a strong, long-lived density peak
survives; a 60-member fluctuation that dissolves as soon as the scale changes
does not, because the cut through the leaf tree is chosen by the trace over all
cuts rather than fixed in advance. That is the behaviour our sample needs, where
M~67 (230 members) and NGC~2158 (6 members) must be judged by the same rule.
What changes is not the units but who chooses: persistence is measured in stars,
and the data pick the value.

\textbf{What does not change.} The cut is still a choice --- PLSCAN makes it
automatically, by maximum total persistence, not without assumptions. The
neighbour count $k$ that sets the core distance is still an input, and their
own sensitivity study is the evidence that this matters less than HDBSCAN$^*$'s
size threshold. The assumption is moved to a quantity whose meaning is more
stable across datasets, and it is made explicit in the barcode, where a user can
see what a different cut would cost.
Nor does any of this repair a bad space: if the field and the clusters are not
separated in the geometry, no persistence criterion will separate them.

\subsection{Where it appears in this workbook}

PLSCAN enters our pipeline through EVoC, and that is worth understanding
before \S\ref{sec:evoc}. EVoC's final step sweeps the clustering of its
internal node embedding across scales, scores each resulting layer by the
total persistence of the clusters it contains, and returns the best-scoring
layer's labels. That is the PLSCAN idea used as a \emph{model-selection rule}
inside a larger algorithm --- the reason EVoC needs no
\texttt{min\_cluster\_size} from the user even though the density clustering
inside it would otherwise want one. On the deck's own terminology: the first
two columns of the EVoC pipeline come from UMAP's graph machinery, the middle
from HDBSCAN$^*$, and the last from PLSCAN.

\begin{issues}[EXERCISES]
\begin{enumerate}
\item Sketch a one-dimensional density with three peaks of different heights
  and widths. For each peak, mark the interval of smoothing scale over which
  it survives as a separate peak, and draw the resulting barcode. Now add a
  fourth peak that is real but tiny --- 2\% of the mass of the smallest other
  peak. Does the barcode distinguish it from a fluctuation?
\item In HDBSCAN$^*$ terms, what is the difference between ``this branch is
  shorter than \texttt{min\_cluster\_size}'' and ``this branch is shorter-lived
  than the others''? Construct a dataset where the two rules give different
  answers, and say which you would trust.
\item The persistence of a cluster is a range of minimum cluster sizes, i.e.\
  a number of stars (their Eq.~7). Using the member counts of
  \textbf{Table~\ref{tab:clusters}}, work out what a persistence of a given
  value implies for M~67 and for NGC~2158, and argue whether the same cut can
  serve both.
\end{enumerate}
\end{issues}
